\section{Conclusão}

As disciplinas de endodontia e prótese dentária encontram nos gêmeos digitais não apenas uma atualização de métodos de imagem e fresagem, mas a redefinição de seus alicerces conceituais, migrando de práticas reativas para uma abordagem estritamente preditiva e matematicamente modelada. Na endodontia, o diagnóstico microestrutural, o planejamento tridimensional do acesso via \textit{Deep Learning} e a simulação de fadiga termo-mecânica de limas de NiTi impulsionam a previsibilidade e mitigações de iatrogenias para níveis outrora inatingíveis \cite{elsaid2025digital, khoshfekr2025review}.

Na prótese dentária e implantodontia, a capacidade in silico de antever a distribuição de estresse, a dissipação de energias na estrutura osso--implante e o ajuste milimétrico dos planos oclusais constituem mecanismos potentes de extensão da longevidade restauradora e maximização do conforto fisiológico e fonético do paciente \cite{roy2025applications, olawade2026advancing}. A materialização de um gêmeo digital unificado que consolide os mundos intra e extrarradicular traduz o paradigma do design \textit{fail-safe}, impedindo a execução de intervenções estruturalmente fadadas à falha \cite{xing2024clinical}.

\destaque{A superação de desafios estruturais -- validados pela necessidade de dados empíricos in vivo, adoção de protocolos não-proprietários e interoperabilidade clínica -- será o divisor de águas na universalização dessa tecnologia.}

Tendo como pano de fundo a infraestrutura de modelagem aberta discutida no Capítulo~\ref{ch:opencode} e as perspectivas de simulação contínua, o gêmeo digital odontológico cessa de ser um recurso puramente virtual para se instituir como o repositório central de saúde oral do paciente. Na esteira dessa evolução, o tratamento torna-se menos dependente da heurística individual do clínico e mais embasado em metadados validados, concretizando a promessa de uma odontologia de precisão plenamente integrada.
